\begin{claim}
Assume that until step $\tau$, there was no deviation cluster of size greater than $d$. Let $\hat{K}$ be a spanned slice induced by a deviation of a cluster state at step $\tau$, and let $\omega$ be an explanation for $\hat{K}$ such that the number of vertices in $\omega$ is greater than $\beta n$. Then there is a subgraph (subtree) $\omega^{\prime}$ of $\omega$ such that $\omega^{\prime}$ has fewer than $\frac{\beta}{2}n$ edges and at least $\frac{\beta}{4\gamma} n$ leaves.
\end{claim}

\begin{proof}
Let $\omega$ be an explanation of $\hat{K}$ with more than $\beta n$ vertices. Since there is no deviation cluster of size greater than $d$ until step $\tau$, we have that every vertex of $\omega$ is a spanned slice induced by a boundary of a deviation cluster of size at most $d$. Thus, the image of $\omega$ over the infinite plane is decomposed into finite assignments, and $\omega$ satisfies the shrink property in \Cref{claim:shrinkslice}, and by \Cref{claim:edgesin}, the ratio of edges to leaves of $\omega$ is at most $\gamma$.
Observe that a spanning tree of $\omega$ has a ratio of edges to leaves that is at most $\omega$'s ratio, and the number of edges equals $\omega$'s vertices minus 1. Therefore, by applying \Cref{claim:big_subtree} on $\omega$'s spanning tree, $\omega$ has a subtree with less than $\frac{\beta}{2}n$ edges and at least $\frac{\beta}{4\gamma} n$ leaves. 
  \end{proof}

 

  \begin{proof}
    Now, since any two connected vertices are at a space-distance of at most $\max\{\Size{\text{arrow}}, \Size{\text{fork}}\} \le d^{\prime}$, it follows that all the leaves in the subtree are at a space-distance of at most $d^{\prime} \beta n$ from each other. Thus, if $\beta < \frac{1}{d^{\prime}}$, no two of them have the same space-coordinate modulo $n$. In particular, each of the leaves has a unique space-coordinate modulo $n$, and therefore the fault subset that induces the explanation contains at least $X$ independent faults. Hence, the probability of having such a subtree is at most $p^{\frac{\beta}{4\gamma}n}$.
  
Because any $\omega$ with more than $\beta n$ vertices implies the existence of a subtree with size less than $\frac{\beta}{2}n$, it's enough to bound the probability of having such a subtree. Considering that the trees are invariant to shifting by $n \cdot \mathbf{1}_{i}$, we have that it's sufficient to show that there are no such trees which are rooted in the finite cube $\mathbb{Z}_{n}^{d} \times \mathbb{Z}_{t}$. Finally, we use \cref{claim:counting_trees} to bound the number of such trees given a fixed root. Overall, we get:
  \begin{equation*}
    \begin{split}
      \prb{\text{No } \omega \text{s with more than  } \beta n \text{ vertices} } \le n^{d} t  \left( \xi^{2}  p^{\frac{\beta}{4\gamma}} \right)^{n}   
    \end{split}
  \end{equation*}
\end{proof}


